\documentclass[a4paper,12pt]{article}
\usepackage{../packages/coursCollege}
\newcommand{\Chapitre}{Dérivation suite}
\renewcommand{\path}{../../}

\renewcommand{\cours}{3MA1~--~EG~--~ns~--~2025-2026}
\begin{document}

\section{Activités}
\begin{activite}{Relation continuité-dérivabilité}
\tcblower
\begin{tasks}(1)
\task Explorer la relation entre continuité et dérivabilité en un point.
\task Énoncer un théorème puis le démontrer.
\task Que penser de la réciproque de ce théorème ?
\end{tasks}
\end{activite}
\begin{activite}{Une limite trigonométrique importante}
\tcblower
\begin{tasks}(1)
\task On considère la fonction $f$ déterminée par $f(x) = \dfrac{\sin(x)}{x}$
\begin{enumerate}
\item Déterminer $D_f$ et $Z_f$.
\item D'après votre intuition, et/ou à l'aide de la calculatrice:
\item Que pensez-vous de $\displaystyle\lim_{x \to 0} f(x)$ et de ? $\displaystyle\lim_{x \to 0} f(x)$
\item Existe-t-il une asymptote horizontale pour cette fonction?
\item Pouvez-vous justifier ces réponses?
\end{enumerate}

\task On s'intéresse maintenant à $\displaystyle\lim_{x \to 0} f(x)$
\begin{enumerate}
\item Conjecturer une valeur pour cette limite, à partir d'essais numériques avec la calculatrice.
\item Proposer une représentation graphique de $f$ sur la base des informations recueillies jusqu'à présent.
\end{enumerate}

\task Que penser de $\displaystyle\lim_{y \to 0} \dfrac{\sin(y)}{y}$, de $\displaystyle\lim_{z \to 0} \dfrac{\sin(z)}{z}$, de $\displaystyle\lim_{x \to a} \dfrac{\sin(x-a)}{x-a}$, de $\displaystyle\lim_{x \to a} \dfrac{\sin(x-a)}{x-a}$ ?

\task Calculer $\displaystyle\lim_{x \to 0} \dfrac{\cos(x)}{x}$ et $\displaystyle\lim_{x \to 0} \dfrac{\tan(x)}{x}$
\end{tasks}
\end{activite}
\begin{activite}{Indéterminations trigonométriques de type 0/0}
\tcblower
\begin{tasks}(1)
\task Calculer :
\begin{multicols}{2}
\begin{enumerate}
\item $\displaystyle\lim_{x \to 0} \dfrac{\sin(12x)}{6x}$
\item $\displaystyle\lim_{x \to 0} \dfrac{\sin(8x)}{10x^2}$
\item $\displaystyle\lim_{x \to 2} \dfrac{\sin(x-2)}{3x-6}$
\item $\displaystyle\lim_{x \to 0} \dfrac{\tan(12x)}{6x}$
\item $\displaystyle\lim_{x \to 2} \dfrac{6-3x}{\sin(x-2)}$
\item $\displaystyle\lim_{x \to 2} \dfrac{\cos(x-2)}{3x-6}$
\end{enumerate}
\end{multicols}
\end{tasks}
\end{activite}
\begin{activite}{Dérivées trigonométriques}
\tcblower
\begin{tasks}(1)
\task Représenter graphiquement la fonction sin et sa dérivée sur un même repère.
\task Quelle conjecture quant à $[\sin(x)]'$ peut-on énoncer ?
\task Démontrer ce résultat.

Indication : utiliser la formule trigonométrique $\sin(a) - \sin(b) = 2\sin\left(\dfrac{a-b}{2}\right)\cos\left(\dfrac{a+b}{2}\right)$

\task Déterminer $[\cos(x)]'$ et $[\tan(x)]'$.
\end{tasks}
\end{activite}
\begin{activite}{Une nouvelle formule de dérivation}
\tcblower
\begin{tasks}(1)
\task Comment déterminer $[\sin(x^3)]'$ ? Quelle est l'opération considérée ici ?
\task Énoncer une nouvelle formule de dérivation.
\task Déterminer :
\begin{multicols}{2}
\begin{enumerate}
\item $[\sin(4x^4-2x)]'$
\item $[\sin^4(x)]'$
\item $[\cos^2(x)]'$
\item $[\tan^4(x^4-1)]'$
\item $[\sin^4(4x^4-2x)]'$
\item $[\sin(\sin(x))]'$
\end{enumerate}
\end{multicols}
\task Quelle relation y a-t-il entre cette nouvelle formule et $\|f(x)\|^n = n[f(x)]^{n-1} \cdot f'(x), \forall n \in \mathbb{R}$
\task Énoncer une formule pour déterminer $\sqrt{f(x)}$.
\end{tasks}
\end{activite}
\begin{activite}{Étude de fonction trigonométrique}
\tcblower
Étudier complètement la fonction $f$ déterminée par $f(x) = \dfrac{1}{2\sin(x)} + 2\sin(x)$.
\end{activite}
\begin{activite}{Démonstration des formules de dérivation}
\tcblower
\begin{tasks}(1)
\task Énoncer les formules de dérivation connues comme des théorèmes en identifiant clairement hypothèses et conclusions.
\task Démontrer les théorèmes pour $\alpha \cdot f$, $f + g$, $f \cdot g$, $\dfrac{1}{f}$.
\task Démontrer le théorème pour $\dfrac{f}{g}$.
\end{tasks}
\end{activite}
\begin{activite}
	Démonstration des formules de dérivation
	\tcblower
	Énoncer et démontrer la formule de dérivation de $x^n$ comme un théorème en identifiant clairement les hypothèses et conclusions.
\end{activite}
\begin{activite}{Image d'un fermé}
\tcblower
On considère le théorème suivant : « Si $f: I \to \mathbb{R}$ est une fonction continue sur $[a,b]$, alors il existe $m$ et $M$ deux nombres réels tels que $f([a,b]) = [m;M]$ »

\begin{tasks}(1)
\task Faire un schéma pour interpréter graphiquement ce résultat.
\task Expliquer pourquoi on l'appelle le théorème « Image d'un fermé par une fonction continue ».
\task Illustrer (représentation graphique, qui sont $a$, $b$, $m$ et $M$ ?) ce théorème dans les cas suivants :
\begin{enumerate}
\item $f:[-1;1] \to \mathbb{R}$ déterminée par $f(x) = x^2$
\item $f:[1;100] \to \mathbb{R}$ déterminée par $f(x) = 2$
\item $f:[0;2] \to \mathbb{R}$ déterminée par $f(x) = x(x-2)^2$
\end{enumerate}
\task Que penser de la conjecture suivante : « Si $f: I \to \mathbb{R}$ est une fonction sur $[a,b]$, alors il existe $m$ et $M$ deux nombres réels tels que $f([a,b]) = [m;M]$ ». Qu'en déduire ?
\end{tasks}
\end{activite}
\begin{activite}{Théorème de Rolle}
\tcblower
On considère le théorème suivant : « Si $f: I \to \mathbb{R}$ est telle que $f$ est continue sur $[a,b]$, $f$ est dérivable sur $]a,b[$ et $f(a) = f(b)$, alors il existe au moins un nombre $c \in ]a,b[$ tel que $f'(c) = 0$ »

\begin{tasks}(1)
\task Faire un schéma pour interpréter graphiquement ce résultat.
\task Illustrer (représentation graphique, qui sont $a$, $b$ et $c$ ?) ce théorème dans les cas suivants :
\begin{enumerate}
\item $f:[-1;1] \to \mathbb{R}$ déterminée par $f(x) = x^2$
\item $f:[1;100] \to \mathbb{R}$ déterminée par $f(x) = 2$.
\item $f:[0;2] \to \mathbb{R}$ déterminée par $f(x) = x(x-2)^2$
\end{enumerate}
\task Discuter la pertinence des hypothèses.
\task * Pourquoi exige-t-on un intervalle ouvert pour la dérivabilité et un intervalle fermé pour la continuité ?
\task * Qui était Rolle ?
\task Démontrer le théorème de Rolle.
\end{tasks}
\end{activite}
\begin{activite}{Théorème des accroissements finis}
\tcblower
On considère le théorème suivant : « Si $f: I \to \mathbb{R}$ est telle que $f$ est continue sur $[a,b]$ et $f$ est dérivable sur $]a,b[$, alors il existe au moins un nombre $c \in ]a,b[$ tel que $f'(c) = \dfrac{f(b) - f(a)}{b - a}$ »

\begin{tasks}(1)
\task Faire un schéma pour interpréter graphiquement ce résultat.
\task Illustrer (représentation graphique, qui sont $a$, $b$ et $c$ ?) ce théorème dans les cas suivants :
\begin{enumerate}
\item $f:[-1;2] \to \mathbb{R}$ déterminée par $f(x) = x^2$
\item $f:[1;100] \to \mathbb{R}$ déterminée par $f(x) = 2$.
\item $f:[0;3] \to \mathbb{R}$ déterminée par $f(x) = x(x-2)^2$
\end{enumerate}
\task Discuter la pertinence des hypothèses.
\task Démontrer ce théorème.
\end{tasks}
\end{activite}
\begin{activite}{Application}
\tcblower
Trouver la ou les valeurs prévues par le théorème des AF pour la fonction $f$ définie par $f(x) = \sqrt{x}$ sur l'intervalle $[25;36]$.
\end{activite}
\begin{activite}{Corollaire du théorème des accroissements finis}
\tcblower
On considère le théorème suivant : « Soit $f$ définie sur un intervalle $I$. Alors on a :
\begin{itemize}
	\item si $f'(x) > 0, \forall x \in I$, alors $f$ est strictement croissante sur $I$
\item si $f'(x) < 0, \forall x \in I$, alors $f$ est strictement décroissante sur $I$
\item si $f'(x) = 0, \forall x \in I$, alors $f$ est constante sur $I$
\end{itemize}
\begin{tasks}(1)
\task Illustrer ce théorème.
\task À quoi est-il utile ?
\task Démontrer ce théorème.
\task Que penser de sa réciproque ?
\end{tasks}
\end{activite}
\section{Dérivées des fonctions trigonométriques}
Théorème [Dérivée du sinus]
\begin{thm}
Dérivées des fonctions trigonométriques
\tcblower
On a 
\begin{itemize}
	\item $[\sin(x)]' = \cos(x)$
	\item $[\cos(x)]' = -\sin(x)$
	\item $[\tan(x)]' = 1 + \tan^2(x) = \dfrac{1}{\cos^2(x)}$
\end{itemize}
\end{thm}


\begin{notation}
	\tcblower
De même que pour $[\sin(x)]^2 = \sin^2(x)$ on note plus simplement $[\sin(x)]' = \sin'(x)$, $[\cos(x)]' = \cos'(x)$ et $[\tan(x)]' = \tan'(x)$.
\end{notation}
\begin{exemple}
	$[\sin(-20x^3 + x)]'$
	\tcblower
On utilise la formule de dérivation de la composition de fonctions. 

$\sin(-20x^3 + x) = g(f(x))$ avec $f(x) = -20x^3 + x$ et $g(x) = \sin(x)$

$[\sin(-20x^3 + x)]' = \cos(-20x^3 + x) \cdot (-20x^3 + x)' = \cos(-20x^3 + x) \cdot (-60x^2 + 1)$


\end{exemple}

\begin{exemple}
	$[\cos^5(x)]'$
	\tcblower
On utilise la formule de dérivation de la composition de fonctions. 

$\cos^5(x) = [\cos(x)]^5 = g(f(x))$ avec $f(x) = \cos(x)$ et $g(x) = x^5$

$[\cos^5(x)]' =$
$[[\cos(x)]^5]' = 5[\cos(x)]^4 \cdot [\cos(x)]' = 5[\cos(x)]^4 \cdot [-\sin(x)] = 5[-\cos(x)]^4 \sin(x)$
\end{exemple}

\begin{exemple}
$[\sin^3(x^2 - 1)]'$
	\tcblower
On utilise la formule de dérivation de la composition de fonctions. 

$\sin^3(x^2 - 1) = [\sin(x^2 - 1)]^3 = g(f(x))$ avec $f(x) = \sin(x^2 - 1)$ et $g(x) = x^5$

$[[\sin(x^2 - 1)]^3]' = 3[\sin(x^2 - 1)]^2 \cdot [\sin(x^2 - 1)]'$

et on doit recommencer pour dériver $[\sin(x^2 - 1)]'$ :

$[\sin(x^2 - 1)]' = \cos(x^2 - 1) \cdot (x^2 - 1)' = \cos(x^2 - 1) \cdot (2x)$

d'où enfin : $[[\sin(x^2 - 1)]^3]' = 3[\sin(x^2 - 1)]^2 \cos(x^2 - 1) \cdot 2x = 6x\sin^2(x^2 - 1)\cos(x^2 - 1)$
\end{exemple}

\subsection{Limites trigonométriques}
\insertexo{1M-byf5j}{false}{exo}{true}
\insertexo{1M-h9w1s}{false}{exo}{true}
\insertexo{1M-jqhm4}{false}{exo}{true}
\insertexo{1M-u68mh}{false}{exo}{true}
\insertexo{1M-areny}{false}{exo}{true}
\insertexo{1M-hhksj}{false}{exo}{true}
\subsection{thm formules}
\insertexo{1M-raguk}{false}{exo}{true}
\subsection{thm importants}
\insertexo{1M-d4c88}{false}{exo}{true}
\insertexo{1M-aef2y}{false}{exo}{true}
\insertexo{1M-n67du}{false}{exo}{true}
\insertexo{1M-5hev4}{false}{exo}{true}
\insertexo{1M-hju27}{false}{exo}{true}
\insertexo{1M-zar8b}{false}{exo}{true}
\insertexo{1M-x5ny2}{false}{exo}{true}
\insertexo{1M-r9x4h}{false}{exo}{true}

\end{document}
